\paragraph{Conjugate-coordinate class: matrix Noether identity and the divergence of the stress (conditional; not closure).}
Work in the Poisson limit with \(F^{ab}=\{X^a,X^b\}\) and \(E_c=-\{X^b,\{X_b,X_c\}\}+\mu^2X_c\), the field equation of \((4g_{\mathrm{YM}}^2)^{-1}\int\mathrm{vol}\,(F_{ab}F^{ab}+2\mu^2X_aX^a)\), for arbitrary functions \(X^a\) on the leaf. The Jacobi identity and the Leibniz rule give \(\{X^a,E_a\}=0\) identically. The tensor \(\mathbb T^{ab}=F^{ac}F^{b}{}_{c}+\mu^2X^aX^b\), which is the variation of the action with respect to the contraction form, satisfies \(\{X_a,\mathbb T^{ab}\}=-F^{b}{}_{c}E^c+\tfrac14\{X^b,F_{cd}F^{cd}\}\); its divergence therefore vanishes on the field equation up to a gradient, and \(\mathbb T^{ab}-\tfrac14\eta^{ab}F_{cd}F^{cd}\) is divergence-free on the field equation exactly, the subtracted term being \(\tfrac14\eta^{ab}F_{cd}F^{cd}\) and not a trace of \(\mathbb T^{ab}\). About \(X^a=T^a\) the linear form obeys \(\{T^a,L(A)_a\}=-(3+\mu^2)\{A^a,T_a\}\), which vanishes exactly for \(\mu^2=-3\). For the massless model the matrix-level identity \([X_a,\mathbb T^{ab}]\propto\) equation of motion is standard; the Poisson-limit form \(\{X^a,E_a\}=0\) and the divergence identity with the \(\mu^2\) term are recorded here as exact identities in the polynomial Poisson algebra, and were verified there for both signs of the bracket. The tensor \(\mathbb T^{ab}\) carries matrix indices and \(\{X_a,\cdot\}\) is a frame derivative, not \(\nabla_\mu\mathbb T^{\mu\nu}\); it differs from the coordinate divergence of the fibre-averaged tensor by terms containing derivatives of the frame coefficients, of relative order \(1/(|q|\ell)\), which are not computed. The paragraph does not identify \(\mathbb T^{ab}\) with the source of an Einstein-type equation.

\paragraph{Conjugate-coordinate class: extra stress and crossover response at the leading symbol (conditional; not closure).}
Keep the setting and notation of the preceding paragraphs at the base points \(y^i=0\), with \(s_i^\mu=S_i^{\mu\rho}q_\rho\), \(Q^{\mu\nu}=\sum_is_i^\mu s_i^\nu\) and \(G_{ij}=s_i\!\cdot\! s_j\) (four-dimensional contraction). Add as inputs a term \(\tfrac12w\int h^{\mu\nu}\mathcal E_{\mu\nu}[h]\), the quadratic form of an induced Einstein--Hilbert term with \(w=-\zeta^2/(16\pi G_{\mathrm{eff}})\) when \(\delta g=\zeta h\), and a conserved source \(T_{\mu\nu}\) coupled as \(\int h^{\mu\nu}T_{\mu\nu}\); the normalisation of the tree-level term and of the coupling is absorbed in \(w\) and in \(\Lambda\), the coefficient of \(\tfrac12c_i\cdot Mc_i\), \(Mc^\nu=q^2c^\nu-q^\nu q\!\cdot\! c\). Since \(s_i\!\cdot\! q=0\), \(\{s_1,s_2,s_3,q\}\) is a basis for \(q^2\neq0\); \(c_i=D_{ij}s_j+\gamma_iq\) with \(\gamma_i\) the gauge part, \(h=H(c)\) with \(H\) the linear map \(c\mapsto h\), which depends only on the symmetric part \(D^{S}\) of \(D\), the three antisymmetric components of \(D\) are the inert amplitudes, \(c^TMc=q^2\mathrm{tr}(DGD^T)\), and \(G=q_\perp^2\mathbf 1-(y^0/\ell)^2\mathbf k\mathbf k^T\) with \(q_\perp^2=q^2+(q_\mu y^\mu)^2/\ell^2\) and \(\mathbf k=\mathbf q/|\mathbf q|\); its eigenvalues are \(g=q_\perp^2\) (twice) and \(\xi=-\sigma q^2\), \(\sigma=(y^4)^2/\ell^2\), along the spatial \(\mathbf q\). The restriction to \(y^i=0\) only provides these closed forms.

\emph{Extra stress.} The amplitude equations are equivalent to \(H^{T}X=-\Lambda Mc\) with \(X=w\mathcal E[h]-T\). In exact rational arithmetic at \((y^0,y^4)=(5,3),(13,5)\) and several rational \(q\) one finds that \(X\) is symmetric (equivalent to the three inert-amplitude equations), that \(q^\mu X_{\mu\nu}=0\), and that \(X_{\mu\nu}=-\tfrac12\Lambda\sum_i\tilde s^i_\mu(Mc_i)_\nu\) with \(\tilde s^i=\sum_kG^{-1}_{ik}s_k\). The amplitude equations determine \(X\) up to \(q_\mu q_\nu\), which \(H^T\) annihilates; \(q^\mu X_{\mu\nu}=0\) follows from the Bianchi identity of \(\mathcal E\) and the conservation of \(T\), which fixes that term, as the ambiguity of the anharmonicity tensor is fixed in~\cite{KumarSteinacker2023,Steinacker2026}. The closed form \(X=-\tfrac12\Lambda\sum_i\tilde s^i\otimes Mc_i=-\Lambda\mathcal L\mathcal E\) expresses the extra stress through the tree-level equations; it is not an independent derivation of its conservation, and the Einstein part \(w\mathcal E\) is the supplied input.

\emph{Crossover.} The operator \(\mathcal L\) is defined on the six-dimensional span of the symmetrised products \(s_is_j\) and is invertible for \(q^2\neq0\) and \(\xi\neq0\). With \(\mathbb L[Z]=(\eta Q)Z+Z(Q\eta)\) one has \(h_T=\mathbb L[Y]\), \(Y=-2X/(\Lambda q^2)\), and eliminating the amplitudes gives exactly \((w+\Lambda\,\mathcal L)\mathcal E=T\), \(\mathcal L=\tfrac12\mathbb L^{-1}\circ(2-P\,\mathrm{tr})\), where \(P\) projects on \(q^\perp\). The operator \(\mathcal L\) has eigenvalues \(1/(2g)\) (two modes), \(1/(g+\xi)\) (two), \((g\pm\sqrt{g(g+8\xi)})/(8g\xi)\) (one each) and zero on the three inert modes. In the first channel \((1+\varrho/(2q_\perp^2))\mathcal E=T/w\) with \(\varrho=\Lambda/w\) the ratio of the tree-level to the induced coefficient, the form \((1-m^2\Box^{-1})\) with \(\Box\to-q_\perp^2\) and \(m^2=\varrho/2\); \(q_\perp^2=k^2+\sigma\omega^2\) is positive definite in \((\omega,\mathbf k)\), not the form \(q^2\). The result is the supplied Einstein part \(w\mathcal E\) plus a non-local, frame-dependent tree-level correction.

\emph{Roots.} \(\det(\Lambda A+wB)\propto\Lambda^3(q^2)^9(\Lambda+2gw)^2(\Lambda+(g+\xi)w)^2(-\Lambda^2+2\Lambda gw+8gw^2\xi)\), which reproduces the generalised eigenvalues of the tree-plus-induced form at the sampled points. Besides \(q^2=0\) there are roots that depend on the direction of \(q\) through the frame: in the first channel \(\omega^2=-(k^2+\varrho/2)/\sigma\), complex for every \(k\) if \(\varrho>0\) and for \(k^2>-\varrho/2\) if \(\varrho<0\); the helicity-zero roots are complex when \(g+8\xi<0\). The relation is elliptic in \(\omega\) rather than wave-like: for \(\varrho>0\) the first channel has no real-frequency root, and for \(\varrho<0\) real roots exist only for \(k^2<|\varrho|/2\), with \(\omega^2<0\) and growth rate proportional to \(k\) above that. The anisotropy is that of the background direction \(\hat y=y/\ell\), through \(q_\perp^2=q^2+(q\!\cdot\!\hat y)^2\), and \(G\) is indefinite for spacelike \(q\); the anisotropy is \(SO(3,1)\)-covariant if \(y^\mu\) is treated as a vector, and is not an explicit breaking of covariance. The sign of \(\varrho\) is that of \(\Lambda/w\); the sign of the induced coefficient is not determined here. These are roots of the leading-symbol dispersion relation, not shown to be poles of a retarded response (constraints and gauge fixing are not treated). The extra modes reported in~\cite{Steinacker2026} are, as extracted, rotation-isotropic with a reduced velocity from the cosmic frame; they belong to the same family of non-boost-invariant effects, and no match is claimed, since the map between the present \(q_\perp,\xi,\sigma\) and the notation there is not established.

This is a statement for the Poisson limit, the stated class with the bilinear term, the leading symbol, tree level plus the supplied induced form and conserved source, and the base points with \(y^i=0\). It derives no value or sign of \(G_{\mathrm{eff}}\) or of \(\varrho\), no masses or subleading terms, no matrix-derived matter coupling, no new derivation of the conservation of the extra stress, and no identification of \(G\) with the emergent metric of Chapter~\ref{ch:vierbein-emergence}; the comparison with the modified Einstein equation of~\cite{KumarSteinacker2023,Steinacker2026} is structural. This changes no ledger label; Problem~F remains open.

